\subsection{The Wilson weight of spatial winding}
\label{cl:sec:row-laws}

Fix one of the four seam matchings. Its active lattice is an ordinary
periodic Wilson lattice or the disjoint union of two such lattices.
Use product probability Haar measure throughout. A straight row is a
complete coordinate circle of links, oriented in the positive direction
and based at coordinate zero. Its holonomy is their ordered product.
All periods below are at least four. We first keep the links outside a
single row fixed, without imposing the cluster restriction $\mathcal G$.

\begin{theorem}[A straight winding row in $\SU(2)$]\label{cl:thm:row-law}
For a row of length $L$, the conditional density $f_L$ of its holonomy
relative to Haar probability satisfies
\[
 \|f_L-1\|_\infty\le\eta_L,\qquad
 \eta_L=\frac{1+u}{(1-u)^3}-1,\quad
 u=\rho^L,\quad \rho=\frac{4b}{1+4b}.
\]
The estimate holds for every exterior link configuration and $b\ge0$.
\end{theorem}
\begin{proof}
No plaquette contains two links from the row. Its links are therefore
conditionally independent. In unit quaternion coordinates each of the
six incident plaquettes contributes a linear functional of norm $b$ to
the exponent of a single link. Its law is a translate of the central
density proportional to $e^{s\cos\theta}$, with $0\le s\le6b$.
Here $\cos\theta$ is the normalized real trace. Integrating one link
at fixed staples is the classical multihit construction \cite{PPR83}.

For the irreducible representation of dimension $\ell+1$, the central
density has Fourier matrix $I_{\ell+1}(s)/I_1(s)$ times the identity
when $s>0$. Indeed the character is
$\sin((\ell+1)\theta)/\sin\theta$ and the Haar angle density is
$(2/\pi)\sin^2\theta$. Multiplication and integration give a numerator
proportional to $I_\ell-I_{\ell+2}=2(\ell+1)I_{\ell+1}/s$.
Dividing by the dimension and the mass gives the formula
\cite{Bum13,Wat44}. Translation multiplies the matrix by a unitary.
At $s=0$ all nonconstant coefficients vanish.

For an integer $\nu\ge1$, let $\mathbb E_\nu$ use the probability
density proportional to $e^{st}(1-t^2)^{\nu-1/2}$ on $[-1,1]$.
The integral for $s^{-\nu}I_\nu(s)$ and its derivative give
$I_{\nu+1}/I_\nu=\mathbb E_\nu t$. Integrating the derivative of
$e^{st}(1-t^2)^{\nu+1/2}$ gives
$(2\nu+1)\mathbb E_\nu t=s\mathbb E_\nu(1-t^2)
\le2s(1-\mathbb E_\nu t)$. Consequently
\[
 0\le I_{\nu+1}(s)/I_\nu(s)
 \le 2s/(2s+2\nu+1)\le\rho .
\]
The product of the $\ell$ successive ratios and the $L$ independent
links gives a Fourier matrix of operator norm at most $\rho^{L\ell}$.
In Fourier inversion its contribution to the density is bounded by
$(\ell+1)^2 u^\ell$. The sum over $\ell\ge1$ converges absolutely and
uniformly and equals $\eta_L$.
\end{proof}

\begin{proposition}[Pinning the holonomies of many rows]\label{cl:prop:row-pins}
Take $k$ edge-disjoint straight rows, each of length $L$. If $\eta_L<1$,
their joint Haar density $F$ satisfies, at every tuple of pins $h$,
\[
 (1-\eta_L)^k\le F(h)\le(1+\eta_L)^k.
\]
Rows may be adjacent. In the original two-sheet ensemble take all rows
in spatial direction one, so $L=n=24T$ and $k=2n^2T$. Denote by $p^I$
the disagreement probability calculated with these holonomies pinned
to the identity in every corner. Then
\[
 \left|\log\frac{p_{\rm diff}}{1-p_{\rm diff}}-
             \log\frac{p^I}{1-p^I}\right|
 \le k\log\frac{1+\eta_n}{1-\eta_n}.
\]
The right hand side tends to zero whenever
$n/(1+4b)-3\log n\longrightarrow+\infty$ with $n\to\infty$.
\end{proposition}
\begin{proof}
When exposing the rows successively, all previously exposed holonomies
are measurable with respect to the exterior of the next row. Conditional
expectation preserves the single-row upper and lower bounds. Multiplying
the resulting conditional densities gives the joint bounds almost
everywhere. To obtain a continuous density, replace the last link in
each row by its holonomy, keeping all preceding links as variables.
Haar invariance gives Jacobian one; the remaining integral of the continuous
Wilson weight is continuous in the pins by compactness. The bounds extend
to every pin because Haar measure has full support.

In each corner this density at $I$ is exactly the pinned partition
integral divided by the ordinary partition integral. The same two
multiplicative bounds hold after summing the weights within either pair
of signs. Taking the ratio of the two sums proves the odds estimate.
Finally $\eta_n=4u+O(u^2)$ as $u\to0$ and
$u\le\exp[-n/(1+4b)]$. Since $k=n^3/12$, the stated growth gives
$k\eta_n\to0$.
\end{proof}

The same estimate gives an absolute probability bound for coherent
central winding. In each of the $k$ rows prescribe either $I$ or $-I$
and require the holonomy to have angle at most $\pi/3$ from that
element. The Haar mass of either cap is
$h=1/3-\sqrt3/(4\pi)$. If $n\ge3(1+4b)$, then $u\le e^{-3}<1/16$
and $1+\eta_n\le(17/16)/(15/16)^3<3/2$. The tower argument bounds
the probability by $[h(1+\eta_n)]^k\le 2^{-k}$ in every corner and
therefore in their Wilson mixture. This estimate concerns straight
lattice holonomies; it does not identify them with renormalized loops.

The pinned partition integral has an explicit action. Regard the active
lattice as a circle of length $n$ times the transverse three-dimensional
graph $\Gamma$ with its square faces. For vertex $y$ of $\Gamma$ take
$g_0(y)=I$ and $g_{j+1}(y)=g_j(y)U_1(j,y)$. The pin gives
$g_n(y)=g_0(y)$. The gauge transformation by $g$ makes every axial link
the identity. Write $V_{j,e}=g_j(y)U_{j,e}g_j(y')^{-1}$ for a transverse
edge $e=(y,y')$. With $x(U)=\ReTr(U)/2$, the remaining integral is
\[
 \int\exp\!\left\{b\sum_{j,p\subset\Gamma}x(P_p(V_j))
       +b\sum_{j,e\subset\Gamma}x(V_{j+1,e}V_{j,e}^{-1})\right\}
       \prod_{j,e}dV_{j,e},\qquad V_n=V_0 .
\]
The first sum ranges over square faces and the second over positive
edges. For fixed axial links these changes preserve the transverse
Haar measures. The $n-1$ free axial links per row then integrate to one.
At equal signs $\Gamma$ consists of two tori of shape $n\times n\times T$;
at opposite signs it is one torus of shape $n\times n\times2T$.
The cross terms couple the $n$ layers at the original value of $b$.
On the growth regime of Proposition~\ref{cl:prop:row-pins}, the lower
bound on the seam odds can thus be studied with all these axial links
removed. A comparison of the two remaining integrals is still required.

One cannot obtain that comparison by discarding the normalizing factors
in the single-link calculation. For a concrete example, keep the
transverse links equal to $I$ along and around a direction-one row.
First set all six parallel neighbouring rows equal to $I$. Alternatively
put three at $g_\delta=\operatorname{diag}(e^{i\delta},e^{-i\delta})$
and three at $g_\delta^{-1}$, with $\cos\delta=1-1/(6b)$ and $b\ge1$.
All other links in the exterior can be set to $I$. The two scalar
staple strengths are $6b$ and $6b-1$. For either exterior every plaquette
formed entirely from exterior links has deficit at most
$1-\cos(2\delta)=O(1/b)$. The normalized row laws both tend uniformly
to Haar when $L/b\to\infty$, $b\to\infty$. Yet the parts of the row
integrals coming from its incident plaquettes have ratio
$[z(6b)/z(6b-1)]^L$, where $z(s)=\int e^{s\cos\theta}\,dg$.
Its logarithm divided by $L$ tends to one: $(\log z)'$ is the tilted
mean of $\cos\theta$, which tends to one because the complement of
any cap at the identity has exponentially smaller weight than a
strictly smaller cap there. Thus near-Haar winding distributions are
compatible with an exponentially large ratio of conditional row weights.

\subsection{Temporal winding and the charged spectrum}
\label{cl:sec:charged-spectrum}
Let $G$ be a compact connected matrix gauge group in its chosen faithful
unitary representation. Give it the bi-invariant metric induced by the
Hilbert--Schmidt norm, and denote its diameter by $D_G$. For $b\ge0$ put
\[
 \delta=(1+6b)^{-1/2},\qquad
 J=\max\{1,\lceil D_G/\delta\rceil\},\qquad
 q_G=1-\frac{1}{8J^2}.
\]
\begin{lemma}[Row contraction for a connected group]\label{cl:lem:general-row}
For a row of length $L\ge4$ the holonomy density satisfies
$\|f_L-1\|_\infty\le e^{12b}q_G^{L-2}$ uniformly in its exterior.
\end{lemma}
\begin{proof}
Write the conditional density of one link as $f=e^U/\int e^U$.
Each of its six plaquette terms has absolute value at most $b$.
For a Lie algebra element $X$ let $R_X$ denote differentiation along
$g\exp(tX)$. The second derivative of a normalized trace term is
bounded by $b\|X\|_{\rm HS}^2$, by unitarity and
$\|X^2\|_1=\|X\|_{\rm HS}^2$ for skew-Hermitian $X$.
Thus $|R_X^2 U|\le6b\|X\|_{\rm HS}^2$. Integrating $R_X^2 f$ over
the group gives zero, so that
\[
 \|R_X\sqrt f\|_2^2=-\tfrac14\int (R_X^2 U) f
 \le\tfrac{6b}{4}\|X\|_{\rm HS}^2 .
\]
For $h=\exp X$ with $\|X\|_{\rm HS}\le\delta$, integrate this derivative
along the path and use Haar invariance. The $L^2$ distance between
$\sqrt{f(\cdot h)}$ and $\sqrt f$ is at most $1/2$. Cauchy--Schwarz
on the difference of their squares bounds total variation by the same
quantity. In particular $\int\min(f(g),f(gh))\,dg\ge1/2$.

In a nontrivial irreducible unitary representation $\pi$, choose unit
vectors $v,w$ realizing $a=\|\int\pi(g)f(g)\,dg\|$ with real matrix
element $a$. Then $\int\|\pi(g)v-w\|^2f(g)\,dg=2(1-a)$.
Integrate the squared triangle inequality for $\pi(gh)v-w$ and
$\pi(g)v-w$ against $\min(f(g),f(gh))$. After translation of the first
term one obtains $\|\pi(h)v-v\|^2\le16(1-a)$. A minimizing geodesic
for the bi-invariant metric writes any group element as at most $J$
such increments. Telescoping gives a squared bound $16J^2(1-a)$
everywhere. Its Haar average on the left equals two, since the Haar
average of $\pi$ is zero. Hence $a\le q_G$.

Convolution by each one-link density has norm at most $q_G$ on
mean-zero $L^2(G)$ by the Peter--Weyl decomposition. Moreover
$\|f-1\|_2^2\le\|f\|_\infty\le e^{12b}$. Apply $L-2$ such contractions
between the first and last centered densities and Cauchy--Schwarz at
the two ends. Conditional independence along the row gives the stated
pointwise bound for their convolution.
\end{proof}

\begin{theorem}[Energy of states outside the gauge-invariant subspace]
\label{cl:thm:charged-gap}
On a finite periodic spatial lattice with periods at least four, let
$K=DCD$ be the one-step Wilson transfer operator in temporal gauge
before projecting onto vertex-gauge invariants. Let $\Pi$ denote the
orthogonal projection onto those invariants. For every $b\ge0$,
\[
 \frac{\|K|_{\ker\Pi}\|}{\|K|_{\operatorname{Ran}\Pi}\|}\le q_G .
\]
For $\SU(2)$ the upper bound can be replaced by $4b/(1+4b)$. Both
bounds are independent of the spatial volume. At lattice spacing $a$,
the energy above the vacuum in $\ker\Pi$ diverges as $a\to0$ whenever
$ab\to0$.
\end{theorem}
\begin{proof}
Use a temporal circle of length $L$ and one temporal winding row per
spatial site. Write $M$ for the number of sites. Simultaneous vertex
gauge transformations commute with $K$. The unrestricted Wilson
integral is $\Tr(\Pi K^L)$, by integrating the temporal links,
as in Proposition~\ref{sc:prop:wilson}. Pin all $M$ temporal winding
holonomies to $I$. The periodic axial gauge above now gives $\Tr K^L$
with the same Haar normalization. Positivity of $K$ follows from its
kinetic kernel: the exponential of the real trace inner product is a
sum of positive tensor-power kernels; multiplication on both sides
by $D$ preserves positivity. On the compact lattice $K$ is
Hilbert--Schmidt and its restriction to invariants is nonzero.
The upper joint-density estimate in Proposition~\ref{cl:prop:row-pins}
uses only the one-row upper bound and the tower argument, which also
apply to Lemma~\ref{cl:lem:general-row}. Consequently
\[
 0\le\frac{\Tr(K^L|_{\ker\Pi})}{\Tr(K^L|_{\operatorname{Ran}\Pi})}
 \le(1+e^{12b}q_G^{L-2})^M-1 .
\]
Fix $b$ and the spatial lattice and let $L\to\infty$. The right hand
side is $O_{b,M}(q_G^L)$. For a positive Hilbert--Schmidt operator
$A$, $\|A\|^L\le\Tr A^L\le\|A\|^{L-2}\Tr A^2$. Taking $L$th roots
proves the volume-independent norm inequality. For $\SU(2)$ use
Theorem~\ref{cl:thm:row-law}, for which $\eta_L=O_b(\rho^L)$ if
$b>0$; at $b=0$ the kinetic operator is constant.

In physical units the general lower energy bound is
$-a^{-1}\log q_G\ge1/(8aJ^2)$, and $J^2\le C_G(1+b)$ for a fixed
group and representation. For $\SU(2)$ and $b>0$ the bound is
$a^{-1}\log(1+1/(4b))\ge1/[a(1+4b)]$. Both diverge under the stated
scaling. The limit in $L$ establishes an inequality for each finite
spatial volume before any continuum limit is taken.
\end{proof}
The physical excitations in the Clay problem are gauge invariant.
They lie in $\operatorname{Ran}\Pi$ orthogonal to the vacuum, and their
energies require a comparison inside that subspace. The charged-sector
estimate isolates the part of the temporal-gauge spectrum removed by
Gauss's law.

\begin{corollary}[Seam odds and gauge-invariant insertions in temporal gauge]
\label{cl:cor:temporal-physical}
For a spatial lattice with $M$ sites write $Z_L=\Tr(\Pi K^L)$,
$Z_L^0=\Tr K^L$ and $\epsilon_L=e^{12b}q_G^{L-2}$, $L\ge4$. Computing
the seam log odds with $Z^0$ in place of $Z$ changes them by at most
$2M\epsilon_T$. For bounded slice operators $F_1,\dots,F_j$ commuting
with $\Pi$ and positive integer time increments summing to $L$, the
two normalized, time-ordered expectations differ by at most
\[
 2\Bigl(\prod_{i=1}^{j}\|F_i\|\Bigr)\bigl[1-(1+\epsilon_L)^{-M}\bigr].
\]
On the cube $n=24T$, at $T\asymp N$, $b=O(\log N)$, both changes
vanish faster than every inverse power of $N$ (for insertions take
$L\asymp N$, $j$ fixed and operator norms growing at most polynomially).
\end{corollary}
\begin{proof}
Set $\delta_L=\log(Z_L^0/Z_L)$. The temporal winding bound gives
$0\le\delta_L\le M\log(1+\epsilon_L)\le M\epsilon_L$. The change in
log odds is $\delta_{2T}-2\delta_T$. Nonnegativity of both terms and
$\epsilon_{2T}\le\epsilon_T$ give the first claim.

In either invariant sector let $A$ be the restriction of $K$ and
restrict each $F_i$ to that sector. For $t_1+\cdots+t_j=L$, Schatten
H\"older with exponents $L/t_i$ gives
\[
 |\Tr(A^{t_1}F_1\cdots A^{t_j}F_j)|
 \le \Bigl(\prod_i\|F_i\|\Bigr)\Tr A^L .
\]
The full normalized expectation is a convex combination of the two
sector expectations with charged weight $1-e^{-\delta_L}$. This proves
the asserted difference bound, including the case of a zero charged
trace. Finally $q_G\le1-c_G/(1+b)$ for a constant $c_G>0$.
Thus $M\epsilon_L\le M\exp[12b-c_G(L-2)/(1+b)]$; under the stated
scaling $M=(24T)^3$ is polynomial and the negative exponent dominates
every multiple of $\log N$. Also $1-(1+\epsilon_L)^{-M}\le M\epsilon_L$.
\end{proof}
This comparison transfers convergence of these correlations between the
two ensembles whenever either converges. It supplies no value for the
physical seam odds and no existence or nontriviality theorem for the
limiting correlations.
